% PCT-SOURCE: pdf=59 print=47 kind=prose
advised to prove some of the following statements directly from the
definitions:

% PCT-SOURCE: pdf=59 print=47 kind=display id=2-49a
\begin{equation*}
  \mathcal{F}[\delta(x)] = \frac{1}{(2\pi)^{n/2}}.
\end{equation*}

% PCT-SOURCE: pdf=59 print=47 kind=display id=2-49
\begin{equation}
  \mathcal{F}(\exp ik\mathbin{\cdot}x) = (2\pi)^{n/2}\delta(p-k)
  \tag{2-49}\label{eq:ch2-fourier-exponential}
\end{equation}

% PCT-SOURCE: pdf=59 print=47 kind=prose
Equations (2-35) and (2-36) are also valid for tempered distributions.

% PCT-SOURCE: pdf=59 print=47 kind=prose
There is one other property of $\mathcal{F}$ which will be needed in the
following: its relation to convolution.  For functions in $\mathcal{S}$,
say, we have

% PCT-SOURCE: pdf=59 print=47 kind=display id=2-50
\begin{align}
  [\mathcal{F}(f*g)](k)
  &= (2\pi)^{-n/2}\int e^{-ik\cdot x}\,dx
       \left[\int f(x-\xi)g(\xi)\,d\xi\right] \notag\\
  &= (2\pi)^{-n/2}\int e^{-ik\cdot\xi}\,d\xi\,g(\xi)
       \left[\int e^{-ik\cdot(x-\xi)}\,d(x-\xi)f(x-\xi)\right] \notag\\
  &= (2\pi)^{n/2}(\mathcal{F}f)(k)(\mathcal{F}g)(k).
  \tag{2-50}\label{eq:ch2-fourier-convolution}
\end{align}

% PCT-SOURCE: pdf=59 print=47 kind=prose
So, apart from a factor $(2\pi)^{n/2}$, $\mathcal{F}$ converts a
convolution into a product.  This can be extended to the convolution of
$f\in\mathcal{S}$ with $T\in\mathcal{S}'$:

% PCT-SOURCE: pdf=59 print=47 kind=display id=2-51
\begin{align}
  [\mathcal{F}(f*T)](g)
  &= (f*T)(\mathcal{F}g)=T(f*\mathcal{F}g) \notag\\
  &= (2\pi)^{n/2}T\{\mathcal{F}[(\overline{\mathcal{F}}f)g]\}
   = (2\pi)^{n/2}(\mathcal{F}T)[(\mathcal{F}f)g] \notag\\
  &= (2\pi)^{n/2}[(\mathcal{F}f)(\mathcal{F}T)](g),
  \tag{2-51}\label{eq:ch2-fourier-distribution-convolution}
\end{align}

% PCT-SOURCE: pdf=59 print=47 kind=prose
where we have used the definitions (2-46) and (2-31) of $\mathcal{F}$ and
$f*T$ and the elementary identity $\overline{\mathcal{F}}\widetilde{f}
=\mathcal{F}f$.  We can write (2-51) as

% PCT-SOURCE: pdf=59 print=47 kind=display id=2-52
\begin{equation}
  f*T=(2\pi)^{n/2}\overline{\mathcal{F}}
       [(\mathcal{F}f)(\mathcal{F}T)],
  \tag{2-52}\label{eq:ch2-fourier-convolution-distribution}
\end{equation}

% PCT-SOURCE: pdf=59 print=47 kind=prose
in which case it displays the first part of

% PCT-SOURCE: pdf=59 print=47 kind=theorem id=2-3
\begin{theorem}[2-3]\label{thm:ch2-fourier-fast-decrease}
The Fourier transform of a tempered distribution of fast decrease is an
infinitely differentiable function which is bounded by a polynomial.
\end{theorem}

% PCT-SOURCE: pdf=59 print=47 kind=proof id=2-3
\begin{proof}
% PCT-SOURCE: pdf=59 print=47 kind=prose
It remains only to prove the polynomial boundedness.  That follows since
every tempered distribution is of the form (2-11).
\end{proof}

% PCT-SOURCE: pdf=59 print=47 kind=heading id=2-3
\subsection{Laplace Transforms and Holomorphic Functions}
\label{sec:ch2-laplace}

% PCT-SOURCE: pdf=59 print=47 kind=prose
It is well known that an integral of the form

% PCT-SOURCE: pdf=59 print=47 kind=display id=2-53
\begin{equation}
  g(x)=\int_0^\infty e^{ikx}f(k)\,dk
  \tag{2-53}\label{eq:ch2-one-sided-laplace}
\end{equation}

% PCT-SOURCE: pdf=59 print=47 kind=prose
where $f$, for concreteness sake, is in $\mathcal{S}(\mathbb{R}^1)$, is a
boundary value of a holomorphic function, namely, the Laplace transform

% PCT-SOURCE: pdf=59 print=47 kind=display id=2-54
\begin{equation}
  g(x+iy)=\int_0^\infty e^{ik(x+iy)}f(k)\,dk
  \tag{2-54}\label{eq:ch2-laplace-example}
\end{equation}

% PCT-SOURCE: pdf=60 print=48 kind=prose
holomorphic in the upper half-plane, $y>0$.  [Clearly, the extra factor
$e^{-ky}$ in (2-54) only serves to improve the convergence of an already
splendidly convergent integral.  For $y>0$,

% PCT-SOURCE: pdf=60 print=48 kind=display
\begin{equation*}
  \frac{dg(z)}{dz}=\lim_{\Delta z\to0}
    \left[\frac{g(z+\Delta z)-g(z)}{\Delta z}\right]
\end{equation*}

% PCT-SOURCE: pdf=60 print=48 kind=prose
exists and is independent of direction.  That is one way to define what
is meant by a holomorphic function of a single complex variable.]

% PCT-SOURCE: pdf=60 print=48 kind=prose
In this section, we present theorems which generalize this example in two
respects.  First, we admit several complex variables and replace the
positive $k$ axis by a convex cone.  Second, we admit distributions rather
than functions as boundary values.  The generalization of the upper
half-plane, $y>0$, turns out to be a so-called tube in which the real parts
of the complex variables are unrestricted but the imaginary parts are
required to be in a cone.  The functions holomorphic in the tube, which
are Laplace transforms of tempered distributions vanishing outside a cone,
are not arbitrary but possess certain characteristic boundedness
properties.  It is essential for applications to field theory to have
these boundedness conditions spelled out in detail.

% PCT-SOURCE: pdf=60 print=48 kind=prose
Since it is only in recent years that holomorphic functions of several
complex variables have turned out to be of use in theoretical physics, we
begin by recalling their definition and basic properties.  We denote by
$\mathbb{C}^n$ the $n$-dimensional vector space with complex scalars.  A
function $f$, defined in a neighborhood of a point $w$ of $\mathbb{C}^n$,
 is said to be \emph{holomorphic}\begingroup
 \renewcommand{\thefootnote}{\ensuremath{\dagger}}
 \footnote{Often also termed \emph{analytic}.}
 \endgroup
at the point $w$ if there exists a multiple power series

% PCT-SOURCE: pdf=60 print=48 kind=display id=2-55
\begin{equation}
  \sum_{k_1,\ldots,k_n=0}^{\infty}
  a_{k_1 k_2\ldots k_n}(z_1-w_1)^{k_1}\cdots
  (z_n-w_n)^{k_n}
  \tag{2-55}\label{eq:ch2-holomorphic-power-series}
\end{equation}

% PCT-SOURCE: pdf=60 print=48 kind=prose
which converges for $z$ in some neighborhood of $w$ and is equal to
$f(z)$ there.  Just as in the case of power series in one variable, one
proves that if (2-55) converges at $z$, it converges absolutely and
uniformly for every $\zeta$ satisfying

% PCT-SOURCE: pdf=60 print=48 kind=display id=2-56
\begin{equation}
  |\zeta_j-w_j|\leq R_j=|z_j-w_j|-\epsilon,
  \qquad j=1,\ldots,n,\qquad\text{for any }\epsilon>0.
  \tag{2-56}\label{eq:ch2-polydisc-radii}
\end{equation}

% PCT-SOURCE: pdf=60 print=48 kind=prose
We shall call such a region a \emph{polydisc}.  There is a theorem of
Weierstrass, familiar from the theory of holomorphic functions of one
complex variable, which says that one can differentiate such a series term
by term.  Using this result one gets

% PCT-SOURCE: pdf=60 print=48 kind=display id=2-57
\begin{equation}
  a_{k_1\ldots k_n}
  =\frac{1}{k_1!\cdots k_n!}
  \left.\frac{\partial^{|k|}f(z_1,\ldots,z_n)}
  {(\partial z_1)^{k_1}\cdots(\partial z_n)^{k_n}}
  \right|_{z_1=w_1,\ldots,z_n=w_n} .
  \tag{2-57}\label{eq:ch2-holomorphic-coefficients}
\end{equation}

% PCT-SOURCE: pdf=61 print=49 kind=prose
Thus, in particular, a holomorphic function of $z_1,\ldots,z_n$ is an
infinitely differentiable function of the $n$ real parts of the $z$'s if we
take the imaginary parts equal to zero.  Just how special it is can be seen
more clearly if one uses Cauchy's integral formula $n$ times to obtain

% PCT-SOURCE: pdf=61 print=49 kind=display id=2-58
\begin{equation}
  f(z_1,\ldots,z_n)=\frac{1}{(2\pi i)^n}
  \oint_{|\zeta_1-w_1|=R_1}\!\cdots\!\oint_{|\zeta_n-w_n|=R_n}
  \frac{f(\zeta_1,\ldots,\zeta_n)\,d\zeta_1\cdots d\zeta_n}
  {(\zeta_1-z_1)\cdots(\zeta_n-z_n)}
  \tag{2-58}\label{eq:ch2-cauchy-polydisc}
\end{equation}

% PCT-SOURCE: pdf=61 print=49 kind=prose
as a representation valid for $z_1,\ldots,z_n$ in a polydisc of uniform
convergence of (2-55).  From this there easily follows the estimate

% PCT-SOURCE: pdf=61 print=49 kind=display id=2-59
\begin{equation}
  \frac{1}{k_1!\cdots k_n!}
  \left|\frac{\partial^{|k|}f(z_1,\ldots,z_n)}
  {(\partial z_1)^{k_1}\cdots(\partial z_n)^{k_n}}\right|
  \leq\frac{C}{R_1^{k_1}\cdots R_n^{k_n}} .
  \tag{2-59}\label{eq:ch2-cauchy-estimate}
\end{equation}

% PCT-SOURCE: pdf=61 print=49 kind=prose
These inequalities are essentially all that distinguishes an infinitely
differentiable function from a holomorphic function.  In no physical
context we know is (2-59) a direct physical statement as it stands, but it
has been shown that in a number of physical situations, notably the ones
considered in Chapters 3 and 4, the physical laws involved imply the
inequalities (2-59) for important functions occurring in the theory.  This
has highly remarkable consequences as we shall see.

% PCT-SOURCE: pdf=61 print=49 kind=prose
There are two important and related principles governing the behavior of
holomorphic functions of a single complex variable which are also valid
for holomorphic functions of several complex variables: \emph{analytic
continuation} and determination of a function by its values on a \emph{real
environment}.  According to the first, one can extend the domain of a
holomorphic function in a unique way by considering sequences of
overlapping polydiscs.  Of course, in general, using two different sequence
of overlapping discs will lead to different values at the same point of
$\mathbb{C}^n$, and one is forced to introduce a Riemann ``surface'' for the
function in order to restore the single valuedness.  This will not be
necessary in any of the applications we have in mind.  The second principle
arises from the fact that all the coefficients of the power series for a
holomorphic function can be obtained by calculating the derivatives in
(2-57) in the direction of the real axis.  Thus, a holomorphic function is
determined in a full complex neighborhood of a point of $\mathbb{C}^n$ by
its values on a real environment, that is, on an open set of $\mathbb{R}^n$
obtained by letting only the real parts of the complex variables vary.

% PCT-SOURCE: pdf=61 print=49 kind=prose
An important example of the unique determination of a holomorphic function
by its values on a real environment occurs in the study of sets of
holomorphic functions with a transformation law under $L^{\uparrow}_{+}$
(see

% PCT-SOURCE: pdf=62 print=50 kind=prose
Theorems 2-11 and 3-5).  There we have a function defined on the restricted
Lorentz group and want to assert that its extension to the complex Lorentz
group is unique.  This is an immediate consequence of what has been said
above, provided we can find a parametrization of the Lorentz group such
that the parameters, say $\lambda_1,\ldots,\lambda_6$, are real and
independent for the real group, and complex and independent for the complex
group, and the function in question is holomorphic in them.  The
parametrization need only work for a neighborhood, $N$, of the identity
element because we can get a parametrization for a neighborhood of an
arbitrary element, $g$, by multiplying all the elements of $N$ by $g$ and
using the same parametrization for the neighborhood $gN$ of $g$.

% PCT-SOURCE: pdf=62 print=50 kind=prose
In the practical applications there will be two kinds of functions:

% PCT-SOURCE: pdf=62 print=50 kind=display id=2-60
\begin{equation}
  F(\Lambda_{s_1},\ldots,\Lambda_{s_n})
  \tag{2-60}\label{eq:ch2-lorentz-holomorphic-function}
\end{equation}

% PCT-SOURCE: pdf=62 print=50 kind=display id=2-60b
\begin{equation*}
  S(\Lambda)_{\alpha\beta},
\end{equation*}

% PCT-SOURCE: pdf=62 print=50 kind=prose
where $F$ is holomorphic in its arguments.  Since a holomorphic function
of a holomorphic function is holomorphic, the first function will be of the
required kind if an appropriate parametrization of $\Lambda$ has been
found.  The matrix $S(\Lambda)$ is one of the irreducible representations
$\mathcal D^{(j/2,k/2)}$ of $SL(2,\mathbb{C})$, with $j+k$ even.  [It is therefore a
representation of $L^{\uparrow}_{+}$ because
$\mathcal D^{(j/2,k/2)}(A)=\mathcal D^{(j/2,k/2)}(-A)$ in that case.]  Now
$\mathcal D^{(j/2,k/2)}(A)_{\alpha\beta}$ is a polynomial in the matrix elements of
$A$, and their complex conjugates and the matrix elements of $A$ can be
expressed locally as analytic functions of the matrix elements of
$\Lambda(A)$, so again parametrizing $\Lambda$ properly leads to the
desired expression.  Incidentally, the analytic continuation is given by
$\mathcal D^{(j/2,k/2)}(A,B)_{\alpha\beta}$, which can be expressed locally as a
function of the matrix elements of $\Lambda(A,B)$.

% PCT-SOURCE: pdf=62 print=50 kind=prose
Thus, it remains to give an appropriate parametrization of $\Lambda$.  For
this purpose, note that the matrix $\Lambda^\mu{}_{\nu}$ can be written as
an exponential of a matrix $\Sigma^\mu{}_{\nu}$, and

% PCT-SOURCE: pdf=62 print=50 kind=display id=2-60c
\begin{align*}
  \Lambda^T G\Lambda=G
  &\quad\text{is equivalent to}\quad
  \Sigma^T=-G\Sigma G \\
  \text{or}\qquad
  \Sigma_{\mu\nu}&=-\Sigma_{\nu\mu} .
\end{align*}

% PCT-SOURCE: pdf=62 print=50 kind=prose
The six parameters $\Sigma_{01}$, $\Sigma_{02}$, $\Sigma_{03}$,
$\Sigma_{12}$, $\Sigma_{13}$, $\Sigma_{23}$ are arbitrary real for
$L^{\uparrow}_{+}$ and arbitrary complex for $L_{+}(\mathbb{C})$, and can
be taken for $\lambda_1,\ldots,\lambda_6$, provided one considers a
sufficiently small neighborhood of the identity.

% PCT-SOURCE: pdf=62 print=50 kind=prose
It is natural to ask whether the notion of real environment is also
significant when it lies on the boundary of the domain in which a function
is holomorphic.  For one complex variable, that is well known to be the
case; a generalization to $n$ variables will be obtained in Theorem 2-17.

% PCT-SOURCE: pdf=63 print=51 kind=prose
There is a useful criterion for a function of several variables to be
holomorphic which we shall use several times in the following.

% PCT-SOURCE: pdf=63 print=51 kind=theorem id=2-4
\begin{theorem}[2-4]\label{thm:ch2-separate-holomorphy}
Let $F$ be a function defined in some open set $D$ of $\mathbb{C}^n$.  A
necessary and sufficient condition that $F$ be holomorphic is that it be
continuous in all variables together and holomorphic in each variable
separately.
\end{theorem}

% PCT-SOURCE: pdf=63 print=51 kind=remark id=2-4
\begin{remark*}
This theorem is also true if the condition of joint continuity is dropped.
That is a deep result known as Hartog's theorem.  (See Bochner and Martin,
Chapter VII.)  The analogue of Hartog's theorem for infinitely
differentiable functions is false; this is another indication of the special
character of holomorphic functions.
\end{remark*}

% PCT-SOURCE: pdf=63 print=51 kind=proof id=2-4
\begin{proof}
% PCT-SOURCE: pdf=63 print=51 kind=prose
In a suitable polydisc of center $w_1,\ldots,w_n$ contained in $D$ we can
write (using the separate holomorphy) the successive integral

% PCT-SOURCE: pdf=63 print=51 kind=display
\begin{equation*}
\begin{aligned}
  F(z_1,\ldots,z_n)
  &=\frac{1}{2\pi i}\oint_{|\zeta_1-w_1|=R_1}
    \frac{d\zeta_1}{\zeta_1-z_1}\frac{1}{2\pi i}
    \oint_{|\zeta_2-w_2|=R_2}\frac{d\zeta_2}{\zeta_2-z_2}\cdots \\
  &\qquad\cdots\frac{1}{2\pi i}
    \oint_{|\zeta_n-w_n|=R_n}\frac{d\zeta_n}{\zeta_n-z_n}
    F(\zeta_1,\ldots,\zeta_n).
\end{aligned}
\end{equation*}

% PCT-SOURCE: pdf=63 print=51 kind=prose
But, because of the joint continuity, the successive integral can be
written as a multiple integral.  Writing the series expansion

% PCT-SOURCE: pdf=63 print=51 kind=display
\begin{equation*}
  \frac{1}{(\zeta_1-z_1)\cdots(\zeta_n-z_n)}
  =\sum_{k_1,\ldots,k_n=0}^{\infty}
  \frac{(z_1-w_1)^{k_1}}{(\zeta_1-w_1)^{k_1+1}}\cdots
  \frac{(z_n-w_n)^{k_n}}{(\zeta_n-w_n)^{k_n+1}},
\end{equation*}

% PCT-SOURCE: pdf=63 print=51 kind=prose
which converges uniformly in every subpolydisc, and interchanging the order
of integration and summation, we get a convergent series expansion for
$F$.
\end{proof}

% PCT-SOURCE: pdf=63 print=51 kind=prose
Like every continuous function, a holomorphic function, $F$, defined on an
open set $\mathcal{O}$ is a distribution in $\mathcal{D}(\mathcal{O})'$:

% PCT-SOURCE: pdf=63 print=51 kind=display
\begin{equation*}
  F(f)=\int dx_1\,dy_1\cdots dx_n\,dy_n\,
  f(x_1,y_1,\ldots,x_n,y_n)F(z_1,\ldots,z_n),
\end{equation*}

% PCT-SOURCE: pdf=64 print=52 kind=prose
where $f$ is any test function with support lying in $\mathcal{O}$.  In
the proof of the edge of the wedge theorem in Section 2-5 we shall need to
know the relation between convergence in $\mathcal{D}(\mathcal{O})'$ and
uniform convergence on compact sets for a sequence of holomorphic
functions, $F_k$, $k=1,2,\ldots$.  The answer is simple: the two notions are
identical.  Clearly, since a sequence of holomorphic $F_k$ converging
uniformly on compact subsets of $\mathcal{O}$ converges to a function, $F$,
holomorphic in $\mathcal{O}$, we have

% PCT-SOURCE: pdf=64 print=52 kind=display id=2-61
\begin{equation}
  F_k(f)\longrightarrow F(f)
  \tag{2-61}\label{eq:ch2-holomorphic-distribution-convergence}
\end{equation}

% PCT-SOURCE: pdf=64 print=52 kind=prose
for $f\in\mathcal{D}(\mathcal{O})$.  Conversely, if
$\lim_{k\to\infty}F_k(f)$ exists for every $f\in\mathcal{D}(\mathcal{O})$,
we can smear the multiple Cauchy formula

% PCT-SOURCE: pdf=64 print=52 kind=display
\begin{equation*}
  F_k(w_1,\ldots,w_n)=\frac{1}{(2\pi)^n}
  \int_0^{2\pi}d\theta_1\cdots\int_0^{2\pi}d\theta_n\,
  F_k(w_1+R_1e^{i\theta_1},\ldots,w_n+R_ne^{i\theta_n})
\end{equation*}

% PCT-SOURCE: pdf=64 print=52 kind=prose
in the radii $R_1,\ldots,R_n$ with an infinitely differentiable function
$g$ with support lying on the product of the positive real axes and
sufficiently near zero, and such that

% PCT-SOURCE: pdf=64 print=52 kind=display
\begin{equation*}
  \int_0^\infty\!\cdots\!\int_0^\infty
  R_1\,dR_1\cdots R_n\,dR_n\,g(R_1,\ldots,R_n)=1.
\end{equation*}

% PCT-SOURCE: pdf=64 print=52 kind=prose
Then

% PCT-SOURCE: pdf=64 print=52 kind=display
\begin{equation*}
  F_k(w_1,\ldots,w_n)=F_k(f_w),
\end{equation*}

% PCT-SOURCE: pdf=64 print=52 kind=prose
where

% PCT-SOURCE: pdf=64 print=52 kind=display
\begin{equation*}
  f_w(x_1,y_1,\ldots,x_n,y_n)
  =(2\pi)^{-n}g(|z_1-w_1|,\ldots,|z_n-w_n|).
\end{equation*}

% PCT-SOURCE: pdf=64 print=52 kind=prose
Now $f_w\in\mathcal{D}(\mathcal{O})$, so convergence in
$\mathcal{D}(\mathcal{O})'$ implies the convergence of
$F_k(w_1,\ldots,w_n)$ at each point $w_1,\ldots,w_n\in\mathcal{O}$.
Furthermore, since, as $w_1,\ldots,w_n$ varies over a sufficiently small
compact set, $f_w$ varies over a bounded set in $\mathcal{D}(\mathcal{O})$,
the convergence is uniform on $K$.  The limit function $F$ is therefore a
holomorphic function.  This completes our general remarks about holomorphic
functions.  We now turn to the definition and properties of the Laplace
transform.

% PCT-SOURCE: pdf=64 print=52 kind=prose
If $T$ is a distribution in $\mathcal{D}'_p$, it may happen that
$e^{-p\cdot\eta}T$, which is certainly a distribution in $\mathcal{D}'_p$,
is also a distribution in $\mathcal{S}'_p$.  In this case, we can define the
Laplace transform of $T$ as the distribution of $\mathcal{S}'_\xi$ given by

% PCT-SOURCE: pdf=64 print=52 kind=display id=2-62
\begin{equation}
  \mathcal{L}(T)=\mathcal{F}(e^{-p\cdot\eta}T).
  \tag{2-62}\label{eq:ch2-laplace-distribution}
\end{equation}

% PCT-SOURCE: pdf=64 print=52 kind=prose
$\mathcal{L}(T)$ actually depends parametrically on $\eta$ but we shall
not indicate that explicitly.  For a function $T$, this definition (2-62)
is

% PCT-SOURCE: pdf=64 print=52 kind=display id=2-63
\begin{equation}
  \mathcal{L}(T)(\xi,\eta)
  =\frac{1}{(2\pi)^{n/2}}
    \int e^{-ip\cdot(\xi-i\eta)}T(p)\,dp .
  \tag{2-63}\label{eq:ch2-laplace-function}
\end{equation}

% PCT-SOURCE: pdf=64 print=52 kind=prose
The convention of signs in the Fourier transform which is used here leads
to $\xi-i\eta$ rather than $\xi+i\eta$ in the exponent.  This results

% PCT-SOURCE: pdf=65 print=53 kind=prose
from our effort to agree with customary physical notation in Chapters 3
and 4.  $T$ need not have any special support properties for this definition
to make sense.  For example, if $T(p)=\exp(-|p|^2)$ the Laplace transform
certainly exists for all $\eta$.  Thus, we are dealing with what is
sometimes called the two-sided Laplace transform, which includes the
one-sided transform (2-53) as a special case.

% PCT-SOURCE: pdf=65 print=53 kind=prose
The first fact about the $\eta$'s for which $\mathcal{L}(T)$ exists is given
by Theorem 2-5.

% PCT-SOURCE: pdf=65 print=53 kind=theorem id=2-5
\begin{theorem}[2-5]\label{thm:ch2-laplace-convexity}
Let $T$ be a distribution of $\mathcal{D}'_p$.  The set of all $\eta$ such
that $e^{-p\cdot\eta}T$ is in $\mathcal{S}'_p$ is convex.
\end{theorem}

% PCT-SOURCE: pdf=65 print=53 kind=proof id=2-5
\begin{proof}
% PCT-SOURCE: pdf=65 print=53 kind=prose
Suppose $e^{-p\cdot\eta'}T$ and $e^{-p\cdot\eta''}T$ belong to
$\mathcal{S}'_p$.  Let $\eta=t\eta'+(1-t)\eta''$ with $0\leq t\leq1$.
Define the infinitely differentiable function $a$ by

% PCT-SOURCE: pdf=65 print=53 kind=display id=2-64
\begin{equation}
  a(p)=\exp(-p\cdot\eta)
  [\exp(-p\cdot\eta')+\exp(-p\cdot\eta'')]^{-1}.
  \tag{2-64}\label{eq:ch2-laplace-convexity-factor}
\end{equation}

% PCT-SOURCE: pdf=65 print=53 kind=prose
Then $a$ is bounded.  [For any two real positive numbers, $c<d$ implies
$c^t<d^t$ and therefore $c<c^td^{1-t}<d$.  Applying this to
$e^{-p\cdot\eta'}$ and $e^{-p\cdot\eta''}$ gives $0<a(p)\leq1$.]
Furthermore, all the partial derivatives of $a$ with respect to $p$ are
bounded.  This follows because they are either linear combinations of
products of the form (2-64) or analogous functions in which $\eta$ is
replaced by $\eta'$ or $\eta''$.  Consequently, the identity

% PCT-SOURCE: pdf=65 print=53 kind=display
\begin{equation*}
  \exp(-p\cdot\eta)T
  =a[\exp(-p\cdot\eta')T]+a[\exp(-p\cdot\eta'')T]
\end{equation*}

% PCT-SOURCE: pdf=65 print=53 kind=prose
displays $\exp(-p\cdot\eta)T$ as the sum of two distributions in
$\mathcal{S}'_p$.
\end{proof}

% PCT-SOURCE: pdf=65 print=53 kind=prose
Theorem 2-5 shows that the definition (2-62) of the Laplace transform
associates with every $T\in\mathcal{D}'_p$ a convex set of $\eta$ (possibly
empty!) for which the Laplace transform exists.  Now we take a given set,
$\Gamma$, of $\eta$'s and study what distributions $T$ can have Laplace
transforms defined for $\eta\in\Gamma$.

% PCT-SOURCE: pdf=65 print=53 kind=theorem id=2-6
\begin{theorem}[2-6]\label{thm:ch2-laplace-holomorphy}
Let $\Gamma$ be a convex open set of $\mathbb{R}^n$ and $T$ a distribution
$\in\mathcal{D}'_p$ such that $e^{-p\cdot\eta}T\in\mathcal{S}'_p$ for all
$\eta\in\Gamma$.  Then $\mathcal{L}(T)$ is a

% PCT-SOURCE: pdf=66 print=54 kind=theorem-continuation id=2-6
holomorphic function of $\xi-i\eta$ for all $\xi-i\eta$ in the tube
$\mathbb{R}^n-i\Gamma$.  $\mathcal{L}(T)$ satisfies the boundedness condition

% PCT-SOURCE: pdf=66 print=54 kind=display id=2-65
\begin{equation}
  |\mathcal{L}(T)(\xi-i\eta)|\leq |P_K(\xi)|
  \tag{2-65}\label{eq:ch2-laplace-polynomial-bound}
\end{equation}

% PCT-SOURCE: pdf=66 print=54 kind=theorem-continuation id=2-6
for some polynomial $P_K$ when $\eta$ varies over any compact subset, $K$,
of $\Gamma$.

% PCT-SOURCE: pdf=66 print=54 kind=theorem-continuation id=2-6
Conversely, every function holomorphic in the tube $\mathbb{R}^n-i\Gamma$,
and satisfying (2-65) for some polynomial $P_K$ for each compact subset,
$K$, of $\Gamma$, is the Laplace transform of a uniquely determined
distribution $T\in\mathcal{D}'_p$, such that $e^{-p\cdot\eta}T\in
\mathcal{S}'_p$ for all $\eta\in\Gamma$.
\end{theorem}

% PCT-SOURCE: pdf=66 print=54 kind=remarks id=2-6
\begin{remark*}
\textit{1.} A tube clearly is a subset of $\mathbb{C}^n$ with a special
property of translation invariance: it is invariant under real translations.
In the mathematical literature the notation is customarily arranged so that
the invariance is under purely imaginary translations instead.  This
evidently is a matter of convention.

% PCT-SOURCE: pdf=66 print=54 kind=remark-item id=2-6-2
\textit{2.} In the case of one complex variable, $\xi-i\eta$, and $\Gamma$,
the cone $\eta>0$, the theorem deals with holomorphic functions in the
lower half-plane.  Then the boundedness restriction (2-65) asserts that the
holomorphic function is bounded by a polynomial $P_{a,b}(\xi)$ in every
horizontal strip of the form $0<a\leq\eta\leq b<\infty$.  Since the
polynomial may vary with $a$ and $b$ we have no control over the behavior of
the holomorphic function for large $\eta$ and small $\eta$.  The following
two examples illustrate this point.
\end{remark*}

% PCT-SOURCE: pdf=66 print=54 kind=display id=2-66
\begin{equation}
  e^{-z^2/2}=\frac{1}{\sqrt{2\pi}}
  \int_{-\infty}^{\infty}dk\,e^{ikz}e^{-k^2/2}
  \tag{2-66}\label{eq:ch2-laplace-example-gaussian}
\end{equation}

% PCT-SOURCE: pdf=66 print=54 kind=display id=2-67
\begin{equation}
  \sqrt{\pi}\,\frac{e^{i/(4z)}}{\sqrt{-iz}}
  =\int_0^\infty dk\,e^{ikz}\frac{\cosh\sqrt{k}}{\sqrt{k}} .
  \tag{2-67}\label{eq:ch2-laplace-example-cosh}
\end{equation}

% PCT-SOURCE: pdf=66 print=54 kind=prose
The first grows as $e^{y^2/2}$ along the imaginary axis as $y\to\infty$,
while the second grows as $y^{-1/2}e^{1/(4y)}$ as $y\to0$.

% PCT-SOURCE: pdf=66 print=54 kind=proof id=2-6
\begin{proof}
% PCT-SOURCE: pdf=66 print=54 kind=prose
Our first step is to show that for $\eta$ varying in certain subsets of
$\Gamma$, $e^{-p\cdot\eta}T$ can be written as a linear combination of
distributions in $\mathcal{S}'_p$, each of which is a product of a function
in $\mathcal{S}$ with a distribution in $\mathcal{S}'_p$; that is,
$e^{-p\cdot\eta}T$ is a tempered distribution of fast decrease.  The
Fourier transform of such a distribution is an infinitely differentiable

% PCT-SOURCE: pdf=67 print=55 kind=prose
function of $\xi$ according to Theorem 2-3.  The particular form will also
guarantee that it is infinitely differentiable in $\xi$ and $\eta$ together.
Then we shall show that this function satisfies the Cauchy--Riemann
equations in $\xi$ and $-\eta$ and is therefore a holomorphic function of
$\xi-i\eta$ according to Theorem 2-4.

% PCT-SOURCE: pdf=67 print=55 kind=prose
Let $\eta^{(1)},\ldots,\eta^{(\ell)}$ be a set of vectors of $\Gamma$ such
that their convex hull, $H$, that is, the cone of vectors
$\sum_{j=1}^{\ell}t_j\eta^{(j)}$ for $\sum_{j=1}^{\ell}t_j=1$, $t_j\geq0$,
$j=1,\ldots,\ell$, has a non-empty interior.  If $\eta$ is a vector of
that interior, then the function

% PCT-SOURCE: pdf=67 print=55 kind=display id=2-68
\begin{equation}
  a(p,\eta)=\exp(-p\cdot\eta)
  \left[\sum_{j=1}^{\ell}\exp(-p\cdot\eta^{(j)})\right]^{-1}
  \tag{2-68}\label{eq:ch2-laplace-convex-combination}
\end{equation}

% PCT-SOURCE: pdf=67 print=55 kind=prose
is bounded in $p$.  [The argument is quite analogous to that for the
boundedness of (2-64): If $a_i>0$ then

% PCT-SOURCE: pdf=67 print=55 kind=display
\begin{equation*}
  \sum_{i=1}^{\ell}t_i\log a_i\leq \sup_i\log a_i
  \qquad\text{for }\sum_{i=1}^{\ell}t_i=1,\quad t_i\geq0,
  \quad i=1,\ldots,\ell.
\end{equation*}

% PCT-SOURCE: pdf=67 print=55 kind=prose
Therefore $|a(p,\eta)|\leq1$.]  Furthermore, the same holds for all the
derivatives with respect to $p$, and uniformly provided $\eta$ stays in
$H$.  The derivatives of $a(p,\eta)$ with respect to $p$ and $\eta$ are
bounded by polynomials in $p$ uniformly for $\eta$ in $H$.  Thus, writing

% PCT-SOURCE: pdf=67 print=55 kind=display id=2-69
\begin{equation}
  \exp(-p\cdot\eta)T
  =\sum_{j=1}^{\ell}a(p,\eta)
  [\exp(-p\cdot\eta^{(j)})T]
  \tag{2-69}\label{eq:ch2-laplace-convex-decomposition}
\end{equation}

% PCT-SOURCE: pdf=67 print=55 kind=prose
we display the left-hand side as a linear combination of elements of
$\mathcal{S}'_p$, namely, the square brackets, times an infinitely
differentiable function of $p$, all of whose derivatives with respect to
$p$ are bounded and whose derivatives with respect to $\eta$ are polynomial
bounded.  This is already enough to show that the Fourier transform of
$\exp(-p\cdot\eta)T$ is a distribution in $\mathcal{S}'_\xi$ which is
differentiable in $\eta$; but to get that it is a function we need a
stronger statement than (2-69).

% PCT-SOURCE: pdf=67 print=55 kind=prose
We assert that for any compact subset, $K$, of $\Gamma$, there is an
$\epsilon>0$ such that

% PCT-SOURCE: pdf=67 print=55 kind=display id=2-70
\begin{equation}
  \exp\bigl(\epsilon\sqrt{1+|p|^2}\bigr)
  \exp(-p\cdot\eta)T=T_1\in\mathcal{S}'_p
  \tag{2-70}\label{eq:ch2-laplace-fast-decrease}
\end{equation}

% PCT-SOURCE: pdf=67 print=55 kind=prose
for all $\eta\in K$.  To prove this we pick an $\eta\in K$.  Then, if
$p$ is a real vector satisfying $|p|\leq\epsilon$, the set $\eta+p$ will
lie in the convex hull of a finite number of vectors of $\Gamma$, for
sufficiently small $\epsilon$.  Thus, for these $\eta+p$'s, one can
construct $a(p,\eta+p)$ as in (2-68).  If we define

% PCT-SOURCE: pdf=67 print=55 kind=display id=2-70b
\begin{equation*}
  b(p,q)=\exp\bigl[\epsilon\sqrt{1+|p|^2}\bigr]a(p,q),
\end{equation*}

% PCT-SOURCE: pdf=68 print=56 kind=prose
for $|q-\eta|\leq\epsilon$ we have

% PCT-SOURCE: pdf=68 print=56 kind=display
\begin{align*}
  |b(p,q)|&\leq\exp[\epsilon(1+|p|)]|a(p,q)| \\
  &\leq\exp(\epsilon)\sup_{|p'|\leq\epsilon}
     \exp(-p\cdot p')|a(p,q)| \\
  &\leq\exp(\epsilon)\sup_{|p'|\leq\epsilon}|a(p,q+p')| .
\end{align*}

% PCT-SOURCE: pdf=68 print=56 kind=prose
Thus $b(p,q)$ is bounded in $p$ for all $q$ sufficiently close to $\eta$.
Furthermore, its derivatives with respect to $p$ and $q$ are bounded by
polynomials in $p$, since they are linear combinations of derivatives of
$\exp(\epsilon\sqrt{1+|p|^2})$ and those of $a(p,q)$.  If we write

% PCT-SOURCE: pdf=68 print=56 kind=display id=2-70c
\begin{equation*}
  \exp\bigl(\epsilon\sqrt{1+|p|^2}\bigr)\exp(-p\cdot\eta)T
  =\sum_{j=1}^{\ell}b(p,\eta)
  [\exp(-p\cdot\eta^{(j)})T]
\end{equation*}

% PCT-SOURCE: pdf=68 print=56 kind=prose
we have (2-70) only for $\eta$ in the neighborhood of a point of $K$.
But, since $K$ is compact it can be covered by a finite number of such
neighborhoods, and therefore the assertion is valid throughout $K$ if the
least of the $\epsilon$ is taken.  The validity of (2-70) assures us that
for each fixed compact, $K$,

% PCT-SOURCE: pdf=68 print=56 kind=display
\begin{equation*}
  \exp(-p\cdot\eta)T
  =\exp\bigl(-\epsilon\sqrt{1+|p|^2}\bigr)T_1,
  \qquad T_1\in\mathcal{S}'_p,
\end{equation*}

% PCT-SOURCE: pdf=68 print=56 kind=prose
where $T_1$ is in $\mathcal{S}'_p$, and varies over a bounded set when
$\eta$ varies over $K$.  Thus, $\exp(-p\cdot\eta)T$ is a tempered
distribution of fast decrease as promised, and $\mathcal{L}(T)(\xi,\eta)$
is an infinitely differentiable function of $\xi$ and $\eta$.  Furthermore,
for $\eta\in K$,

% PCT-SOURCE: pdf=68 print=56 kind=display id=2-71
\begin{equation}
  |\mathcal{L}(T)(\xi,\eta)|< |P_K(\xi)|,
  \tag{2-71}\label{eq:ch2-laplace-bound-proof}
\end{equation}

% PCT-SOURCE: pdf=68 print=56 kind=prose
where $P_K$ is some polynomial.  Now

% PCT-SOURCE: pdf=68 print=56 kind=display
\begin{align*}
  \frac{\partial}{\partial\xi_j}\mathcal{L}(T)(\xi,\eta)
  &=\mathcal{F}(-ip_j e^{-p\cdot\eta}T) \\
  &=i\frac{\partial}{\partial\eta_j}\mathcal{L}(T)(\xi,\eta),
\end{align*}

% PCT-SOURCE: pdf=68 print=56 kind=prose
which are the Cauchy--Riemann equations for the complex variables
$\xi_j-i\eta_j$, $j=1,\ldots,n$.  Therefore $\mathcal{L}(T)$ is a
holomorphic function of $\xi-i\eta$.  This completes the proof of the
first half of the theorem.

% PCT-SOURCE: pdf=68 print=56 kind=prose
Conversely, suppose $F(\xi-i\eta)$ is a holomorphic function in the tube
$\mathbb{R}^n-i\Gamma$, where $\Gamma$ is an open convex set, and suppose
that for each compact subset $K$ of $\Gamma$ there exists a polynomial
$P_K$ such that

% PCT-SOURCE: pdf=68 print=56 kind=display id=2-72
\begin{equation}
  |F(\xi-i\eta)|\leq |P_K(\xi)|,\qquad \eta\in K.
  \tag{2-72}\label{eq:ch2-laplace-converse-bound}
\end{equation}

% PCT-SOURCE: pdf=69 print=57 kind=prose
Equation (2-72) permits us to say that, for each $\eta\in\Gamma$,
$F(\xi-i\eta)$ is a distribution in $\xi$ so we can immediately write
down the distribution, $\widehat{F}\in\mathcal{S}'_p$, which, when
multiplied by $e^{p\cdot\eta}$, ought to be that of which $F$ is the
Laplace transform

% PCT-SOURCE: pdf=69 print=57 kind=display
\begin{equation*}
  \widehat{F}(p,\eta)=\overline{\mathcal{F}}_{\xi}
  [F(\xi-i\eta)].
\end{equation*}

% PCT-SOURCE: pdf=69 print=57 kind=prose
If we multiply it by $e^{p\cdot\eta}$ we get a distribution which may no
longer be in $\mathcal{S}'_p$, but must be in $\mathcal{D}'_p$.  If we could
show that

% PCT-SOURCE: pdf=69 print=57 kind=display id=2-73
\begin{equation}
  \frac{\partial}{\partial\eta_j}
  [e^{p\cdot\eta}\widehat{F}(p,\eta)]=0
  \tag{2-73}\label{eq:ch2-laplace-eta-independence}
\end{equation}

% PCT-SOURCE: pdf=69 print=57 kind=prose
we would write

% PCT-SOURCE: pdf=69 print=57 kind=display id=2-74
\begin{equation}
  \widehat{F}(p,\eta)=e^{-p\cdot\eta}T(p),
  \qquad T\in\mathcal{D}'_p
  \tag{2-74}\label{eq:ch2-laplace-inverse-form}
\end{equation}

% PCT-SOURCE: pdf=69 print=57 kind=prose
and the proof would be complete.  However, at the present stage we do not
even know that $\widehat{F}(p,\eta)$ is differentiable in $\eta$, let alone
that (2-73) holds.  To get the differentiability we show that the
holomorphy of $F$ and the inequality (2-72) imply analogous inequalities
for the derivatives of $F$.

% PCT-SOURCE: pdf=69 print=57 kind=prose
We first explain the idea of the proof for the case of one variable.  Then
we take as the compact set inside the tube a closed interval
$0<a\leq\eta\leq b$.  (Its position in the lower half-plane of
$\xi-i\eta=z$ can always be arranged by suitable location of the origin of
coordinates.)  We choose a sufficiently high power, $z^k$, of $z$ so that
$f(z)z^{-k}$ is bounded in the strip $a\leq\eta\leq b$.  Then, with the
contour $C$, shown in Figure 2-1, we obtain the integral representations
for $a<\operatorname{Im}z<b$

% PCT-SOURCE: pdf=69 print=57 kind=display id=2-75
\begin{align}
  \frac{f(z)}{z^{k+2}}
  &=\frac{1}{2\pi i}\oint_C
    \frac{f(\zeta)\,d\zeta}{\zeta^{k+2}(\zeta-z)}, \notag\\
  \frac{df/dz}{z^{k+2}}
  &=(k+2)\frac{f(z)}{z^{k+3}}
    +\frac{1}{2\pi i}\oint_C
    \frac{f(\zeta)\,d\zeta}{\zeta^{k+2}(\zeta-z)^2} .
  \tag{2-75}\label{eq:ch2-cauchy-strip-derivative}
\end{align}

% PCT-SOURCE: pdf=69 print=57 kind=prose
From these representations one immediately derives that $|f(z)|\leq
c|z|^k$ in the strip implies $|df/dz|\leq d|z|^{k+2}$.  For $n$
variables the argument is similar.  One takes a compact subset of $\Gamma$
which is a product of $n$ intervals and argues that
$F(z)(z_1^{-k_1}\cdots z_n^{-k_n})$ will be bounded for sufficiently large
$k_1,\ldots,k_n$ by virtue of (2-72).  (Again, by suitable choice of
origin, the intervals have been taken in the lower half-plane.)  Instead of
(2-75) one has a multiple Cauchy integral formula, but the result is the
same: the polynomial boundedness of the partial derivatives of
$F(\xi-i\eta)$.

% PCT-SOURCE: pdf=070 print=58 kind=figure id=2-1
% Native reconstruction from the rendered source page.  The contour runs
% along three horizontal lines; the arrowheads carry their orientation.
\begingroup
\renewcommand{\thefigure}{2-1}
\renewcommand{\figurename}{FIGURE}
\begin{figure}[t]
  \centering
  \begin{tikzpicture}[>=Stealth,line cap=round,line join=round,
    contour/.style={line width=.75pt}]
    % The imaginary axis of the z plane.
    \draw[->,line width=.7pt] (0,-3.03) -- (0,3.28);
    \node at (1.94,3.15) {$z$};

    % The upper line, traversed to the right.
    \draw[contour,->] (-4.06,1.65) -- (3.94,1.65);

    % The middle line, traversed to the left.
    \draw[contour,->] (4.14,0) -- (-1.11,0);
    \draw[contour] (-1.11,0) -- (-4.06,0);

    % The lower line, traversed to the right.
    \draw[contour,->] (-4.06,-1.60) -- (0.86,-1.60);
    \draw[contour] (0.86,-1.60) -- (4.14,-1.60);

    \node at ( 0.33, 0.31) {$a$};
    \node at (-0.29,-1.92) {$b$};
  \end{tikzpicture}
  \caption{The contour $C$ in the plane of $z=\xi-i\eta$.}
  \label{fig:2-1}
  \label{fig:ch2-contour}
\end{figure}
\endgroup


% PCT-SOURCE: pdf=70 print=58 kind=prose
Since that in turn implies the differentiability of $\widehat{F}$, we can
carry out the differentiations in (2-73):

% PCT-SOURCE: pdf=70 print=58 kind=display
\begin{align*}
  \frac{\partial}{\partial\eta_j}
  [e^{p\cdot\eta}\widehat{F}(p,\eta)]
  &=p_j e^{p\cdot\eta}\widehat{F}(p,\eta)
    +e^{p\cdot\eta}\frac{\partial\widehat{F}}{\partial\eta_j}(p,\eta) \\
  &=e^{p\cdot\eta}\mathcal{F}
    \left[i\frac{\partial}{\partial\xi_j}F(\xi-i\eta)
    +\frac{\partial F}{\partial\eta_j}(\xi-i\eta)\right]=0 .
\end{align*}

\end{proof}

% PCT-SOURCE: pdf=70 print=58 kind=prose
We now turn to the further restrictions on the holomorphic functions
$\mathcal{L}(T)$ which result from the requirements that $T$ have support
in a cone and that $T$ be tempered.

% PCT-SOURCE: pdf=70 print=58 kind=prose
First consider the matter heuristically.  If the support of $T$ lies in a
half-space $p\cdot a>A$, where $a$ is some fixed vector, then, for
$\eta\in\Gamma$, the ``integral''

% PCT-SOURCE: pdf=70 print=58 kind=display id=2-76
\begin{equation}
  \int e^{-ip\cdot(\xi-i\eta)}T(p)\,dp
  \tag{2-76}\label{eq:ch2-laplace-support-integral}
\end{equation}

% PCT-SOURCE: pdf=70 print=58 kind=prose
converges even better if the exponential is replaced by

% PCT-SOURCE: pdf=70 print=58 kind=display
\begin{equation*}
  e^{-ip\cdot(\xi-i\eta)}e^{-t(p\cdot a-B)}
  \qquad\text{for }B<A,\quad t\geq0,
\end{equation*}

% PCT-SOURCE: pdf=70 print=58 kind=prose
i.e.,

% PCT-SOURCE: pdf=70 print=58 kind=display
\begin{equation*}
  e^{tB}\int e^{-ip\cdot[\xi-i(\eta+ta)]}T(p)\,dp
\end{equation*}

% PCT-SOURCE: pdf=70 print=58 kind=prose
ought to exist for all $B<A$ and $t\geq0$.  But that means that if
$\eta\in\Gamma$ then $\eta+ta\in\Gamma$ for all $t\geq0$ when $a$ is any
vector such that the support of $T$ lies in the half-space
$p\cdot a>A$.  Of course, this statement is empty unless $\Gamma$ is
non-empty.  If $T$ is tempered, then $\Gamma$ contains at least one point
$\eta=0$, because for $\eta=0$, the Laplace transform is the Fourier
transform.

% PCT-SOURCE: pdf=71 print=59 kind=prose
It need not have any other point; consider, for example,
$T(p)=(1+|p|^2)^{-2n}$.  However, if $T$ is tempered \emph{and} the support
of $T$ lies in a half-space, the combination of the preceding heuristic
discussion and the fact that $\Gamma$ must be convex forces $\Gamma$ to be
a cone.  That is the situation met in quantum field theory as we shall see.

% PCT-SOURCE: pdf=71 print=59 kind=theorem id=2-7
\begin{theorem}[2-7]\label{thm:ch2-laplace-halfspace}
Suppose $T\in\mathcal{D}'_p$ and $\Gamma$ (a convex set of $\mathbb{R}^n$)
consists of vectors $\eta$ such that $e^{-p\cdot\eta}T\in\mathcal{S}'_p$.
If $T$ has its support lying in a half-space $p\cdot a>A$, then $\Gamma$
contains all points of the form $\eta+ta$ with $\eta\in\Gamma$ and
$t\geq0$.
\end{theorem}

% PCT-SOURCE: pdf=71 print=59 kind=proof id=2-7
\begin{proof}
% PCT-SOURCE: pdf=71 print=59 kind=prose
Let $\epsilon>0$ and choose an infinitely differentiable function $q$ of
one variable so that it is 1 for $x\geq A$ and 0 for $x\leq A-\epsilon$.
Define $q_\epsilon(p)=q(p\cdot a)$.  The distribution
$\exp[-p\cdot(\eta+ta)+At]T$ is in $\mathcal{D}'_p$ for $\eta\in\Gamma$,
and has support in the half-space $p\cdot a\geq A$ because $T$ does.
Furthermore, if $f\in\mathcal{D}$, $t\geq0$,

% PCT-SOURCE: pdf=71 print=59 kind=display id=2-77
\begin{equation}
  \{\exp[-p\cdot(\eta+ta)+At]T\}(f)
  =[\exp(-p\cdot\eta)T]
   \{\exp(-tp\cdot a+At)q_\epsilon f\}
  \tag{2-77}\label{eq:ch2-laplace-halfspace-test}
\end{equation}

% PCT-SOURCE: pdf=71 print=59 kind=prose
is independent of $q_\epsilon$ provided it satisfies the above requirements
for some $\epsilon$.  Now the right-hand side of (2-77) has an extension to
$\mathcal{S}'_p$ by continuity because
$[\exp(-tp\cdot a+At)q_\epsilon]f$ is in $\mathcal{S}$ if $f$ is, and
depends continuously on $f$.  Therefore,
$\exp[-p\cdot(\eta+ta)]T\in\mathcal{S}'_p$, so $\eta+ta\in\Gamma$.
\end{proof}

% PCT-SOURCE: pdf=71 print=59 kind=prose
In the applications of Chapters 3 and 4, $T$ is defined for arguments
$p_1,\ldots,p_n$ where each $p_j$ is a real four-vector, and vanishes
unless each $p_j$ lies in or on the plus light cone.  This phrase will
occur so frequently in the following that we introduce the notation
$V_+$ to stand for all real four-vectors $p$ satisfying

% PCT-SOURCE: pdf=71 print=59 kind=display id=2-notation-cone
\begin{equation*}
  p^2=-(p^0)^2+\mathbf{p}^{\,2}<0,\qquad p^0>0,
\end{equation*}

% PCT-SOURCE: pdf=71 print=59 kind=prose
and $\overline{V}_+$ to stand for the closure of $V_+$, the set of $p$
satisfying $p^2\leq0$, $p^0\geq0$.  Since, in Chapters 3 and 4, $T$ will
also be tempered, application of Theorem 2-7 yields that
$\mathcal{L}(T)(\xi-ia)$ is analytic for all $a$ of the form
$a=(a_1,\ldots,a_n)$ with $a_j\in V_+$, $j=1,\ldots,n$.  In the following,
the tube $\mathbb{R}^n-i\Gamma$, where $\Gamma=(a_1,\ldots,a_n)$ with
$a_j\in V_+$, $j=1,\ldots,n$, will be denoted $\mathcal{T}_n$, and
sometimes called the tube without further specification.

% PCT-SOURCE: pdf=71 print=59 kind=prose
The Laplace transforms of distributions with support in a half-space
possess stronger boundedness properties than (2-72).  While these can

% PCT-SOURCE: pdf=72 print=60 kind=prose-continuation
be proved for general tubes, we get an especially neat statement for
$\mathcal{T}_n$ so we confine our attention to that case.

% PCT-SOURCE: pdf=72 print=60 kind=theorem id=2-8
\begin{theorem}[2-8]\label{thm:ch2-laplace-forward-cone}
Let $T\in\mathcal{D}'_p$ and $e^{-p\cdot\eta}T\in\mathcal{S}'_p$ for all
$\eta\in\Gamma$.  Here $p\cdot\eta$ stands for
$\sum_{j=1}^n\sum_{\mu=0}^3p_{j\mu}\eta_j^\mu$ and $\Gamma$ for the cone
$\eta_j\in V_+$, $j=1,\ldots,n$.  Suppose $p\in\operatorname{supp}T$ implies
$p_j\in\overline{V}_+$, $j=1,\ldots,n$.  Then, for each $\eta\in\Gamma$,
there is a polynomial $P_\eta$ such that

% PCT-SOURCE: pdf=72 print=60 kind=display id=2-78
\begin{equation}
  |\mathcal{L}(T)[\xi-i(\eta+a)]|
  \leq |P_\eta(\xi-ia)|
  \tag{2-78}\label{eq:ch2-laplace-forward-cone-bound}
\end{equation}

% PCT-SOURCE: pdf=72 print=60 kind=theorem-continuation id=2-8
for all $\xi$ and all $a\in\Gamma$.

For the converse, if $F$ is a function holomorphic in
$\mathcal{T}_n=\mathbb{R}^{4n}-i\Gamma$ and satisfying the inequality (2-78)
for each $\eta\in\Gamma$ and some polynomial $P_\eta$, then $F$ is the
Laplace transform of a distribution with support in $\Gamma$.
\end{theorem}

% PCT-SOURCE: pdf=72 print=60 kind=proof id=2-8
\begin{proof}
% PCT-SOURCE: pdf=72 print=60 kind=prose
The proof of (2-78) uses the Bros--Epstein--Glaser lemma.  See pages 21--27
in M. Reed and B. Simon, \emph{Methods of Modern Mathematical Physics II,
Fourier Analysis, Self-Adjointness}, Academic Press, 1975.

% PCT-SOURCE: pdf=72 print=60 kind=prose
Conversely, suppose $F$ is holomorphic in $\mathbb{R}^{4n}-i\Gamma=\mathcal{T}_n$
and is bounded according to (2-78).  Then it certainly satisfies (2-72), so
it is the Laplace transform of a distribution $T\in\mathcal{D}'_p$ such
that $e^{-p\cdot\eta}T\in\mathcal{S}'_p$ for all $\eta\in\Gamma$.  It
remains to show that the support of $T$ lies inside $\Gamma$.  Pick a test
function $g$ of compact support lying entirely outside $\Gamma$.  Then

% PCT-SOURCE: pdf=72 print=60 kind=display
\begin{align*}
  T(g)
  &=[e^{-p\cdot(\eta+a)}T][e^{p\cdot(\eta+a)}g] \\
  &=\{\mathcal{F}[e^{-p\cdot(\eta+a)}T]\}
    \{\overline{\mathcal{F}}[e^{p\cdot(\eta+a)}g]\} \\
  &=\int F[\xi-i(\eta+a)]G[\xi-i(\eta+a)]\,d\xi,
\end{align*}

% PCT-SOURCE: pdf=72 print=60 kind=prose
where

% PCT-SOURCE: pdf=72 print=60 kind=display
\begin{equation*}
  G(\xi-i\eta)=(2\pi)^{-2n}
  \int e^{ip\cdot(\xi-i\eta)}g(p)\,dp .
\end{equation*}

% PCT-SOURCE: pdf=72 print=60 kind=prose
Now, by appropriate choices of $a$ it can be arranged that, throughout the
support of $g$, we have $p\cdot(\eta+a)<0$, and as $a\to\infty$,
$p\cdot(\eta+a)\to-\infty$.  Then, as $a\to\infty$ in $\Gamma$,
$G[\xi-i(\eta+a)]\to0$ in $\mathcal{S}_\xi$.  Therefore

% PCT-SOURCE: pdf=72 print=60 kind=display
\begin{align*}
  |T(g)|
  &\leq\int d\xi\,|P_\eta(\xi-ia)|
    |G[\xi-i(\eta+a)]| \\
  &\leq\int\frac{dp}{[1+|p|^2]^k}
    \sup_\xi[1+|\xi|^2]^k
    |P_\eta(\xi-ia)G[\xi-i(\eta+a)]|\longrightarrow0
\end{align*}

% PCT-SOURCE: pdf=72 print=60 kind=prose
as $a\to\infty$ for sufficiently large fixed $k$.
\end{proof}

% PCT-SOURCE: pdf=73 print=61 kind=prose
It is already visible in the example (2-67) that a holomorphic function
which is a Laplace transform need not have a boundary value even in the
sense of distribution theory.  If it is the Laplace transform of a
tempered distribution, $T$, however, the boundary value is just the
Fourier transform of $T$.

% PCT-SOURCE: pdf=73 print=61 kind=theorem id=2-9
\begin{theorem}[2-9]\label{thm:ch2-laplace-boundary-value}
Suppose $T\in\mathcal{S}'_p$ and $\mathcal{L}(T)$ exists for all
$\eta\in\Gamma$ as described in Theorem 2-8.  Then

% PCT-SOURCE: pdf=73 print=61 kind=display id=2-79
\begin{equation}
  \lim_{\eta\to0}\int\mathcal{L}(T)(\xi-i\eta)f(\xi)\,d\xi
  =[\mathcal{F}(T)](f)
  \tag{2-79}\label{eq:ch2-laplace-boundary}
\end{equation}

% PCT-SOURCE: pdf=73 print=61 kind=theorem-continuation id=2-9
that is, $\mathcal{L}(T)$ converges in $\mathcal{S}'_\xi$ to
$\mathcal{F}(T)$ as $\eta\to0$ inside any closed cone in $\Gamma$.

Conversely, if $\mathcal{L}(T)$ converges in $\mathcal{S}'_\xi$ as
$\eta\to0$ in any such cone, $T$ is a tempered distribution.
\end{theorem}

% PCT-SOURCE: pdf=73 print=61 kind=proof id=2-9
\begin{proof}
% PCT-SOURCE: pdf=73 print=61 kind=prose
The first statement of the theorem follows directly from (2-69) and the
argument given in connection with it; $\exp(-p\cdot\eta)T$ is a continuous
function of $\eta$ with values in $\mathcal{S}'_p$ in the convex hull of any
finite set of vectors of $\Gamma$.

% PCT-SOURCE: pdf=73 print=61 kind=prose
The converse statement means that

% PCT-SOURCE: pdf=73 print=61 kind=display
\begin{equation*}
  \left[\lim_{\eta\to0}\mathcal{F}(e^{-p\cdot\eta}T)\right](f)
  =\lim_{\eta\to0}T(e^{-p\cdot\eta}\mathcal{F}f)
\end{equation*}

% PCT-SOURCE: pdf=73 print=61 kind=prose
exists for every $f\in\mathcal{S}$.  The completeness of $\mathcal{S}'$
implies there exists a tempered distribution $T_1$ such that this limit is
$T_1(\mathcal{F}f)$.  But for $\mathcal{F}f\in\mathcal{D}$, clearly
$T_1(\mathcal{F}f)=T(\mathcal{F}f)$.  Therefore $T_1=T$ and $T$ is
tempered.
\end{proof}

% PCT-SOURCE: pdf=73 print=61 kind=prose
Theorem 2-9 gives a characterization of those holomorphic functions which
arise by Laplace transformation from tempered distributions but it is not
expressed very directly in terms of $\mathcal{L}(T)$ as a holomorphic
function.  The last theorem of this section gives a more direct statement.

% PCT-SOURCE: pdf=73 print=61 kind=theorem id=2-10
\begin{theorem}[2-10]\label{thm:ch2-laplace-tempered-bound}
Let $e^{-p\cdot\eta}T\in\mathcal{S}'_p$ for $\eta=0$ and $\eta\in\Gamma$,
with $\Gamma$ as described in Theorem 2-8.  Then, for each compact subset
$K$ of $\Gamma$, there is a polynomial $P_K$ and an integer $r$ such that

% PCT-SOURCE: pdf=73 print=61 kind=display id=2-80
\begin{equation}
  |\mathcal{L}(T)(\xi-it\eta)|\leq\frac{P_K(\xi)}{t^r}
  \tag{2-80}\label{eq:ch2-laplace-tempered-bound}
\end{equation}

% PCT-SOURCE: pdf=74 print=62 kind=theorem-continuation id=2-10
for all $\xi$, all $t$ with $0<t<1$, and all $\eta\in K$.

Conversely, if $F$ is a function holomorphic in $\mathcal{T}_n$ and
satisfying (2-80), $F=\mathcal{L}(T)$, where $T$ is a tempered
distribution.
\end{theorem}

% PCT-SOURCE: pdf=74 print=62 kind=proof id=2-10
\begin{proof}
% PCT-SOURCE: pdf=74 print=62 kind=prose
To get (2-80), we repeat part of the argument used in the proof of
Theorem 2-6.  It shows that we can write
$e^{-tp\cdot\eta}T=a(p,t\eta)T_1$, where $T_1\in\mathcal{S}'_p$ and
$a\in\mathcal{S}_p$ for all $\eta\in K$, $0<t<1$ and $a$ varies over a
bounded set as $\eta$ varies.  Therefore, $\mathcal{L}(T)(\xi-it\eta)$ is
the tempered distribution $T_1$ evaluated for the test function
$(2\pi)^{-2n}a(p,t\eta)e^{-ip\cdot\xi}$.  Consequently, for some integers
$k,l$ and some constant $C$

% PCT-SOURCE: pdf=74 print=62 kind=display id=2-81
\begin{equation}
  |\mathcal{L}(T)(\xi-it\eta)|
  \leq C\left\|e^{-ip\cdot\xi}a(p,t\eta)\right\|_{k,l}
  \tag{2-81}\label{eq:ch2-laplace-tempered-seminorm}
\end{equation}

% PCT-SOURCE: pdf=74 print=62 kind=prose
A straightforward estimate shows that the right-hand side is $\leq
t^{-r}P_K(\xi)$ for some integer $r$ as $\eta$ varies over $K$.

% PCT-SOURCE: pdf=74 print=62 kind=prose
To verify the converse statement, we choose $f\in\mathcal{S}$ and study

% PCT-SOURCE: pdf=74 print=62 kind=display
\begin{equation*}
  \int F(\xi-it\eta)f(\xi)\,d\xi=h(t)
\end{equation*}

% PCT-SOURCE: pdf=74 print=62 kind=prose
as a function of $t$.  We have

% PCT-SOURCE: pdf=74 print=62 kind=display
\begin{align*}
  \frac{dh}{dt}(t)
  &=\int\sum_j\frac{\partial}{\partial(\xi-it\eta)_j}
    F(\xi-it\eta)(-i\eta_j)f(\xi)\,d\xi \\
  &=\int F(\xi-it\eta)i\eta\mathbin{\cdot}
    \frac{\partial}{\partial\xi}f(\xi)\,d\xi \\
  &\ \vdots \\
  h^{(j)}(t)
  &=\int F(\xi-it\eta)
    \left(i\eta\mathbin{\cdot}\frac{\partial}{\partial\xi}\right)^j
    f(\xi)\,d\xi .
\end{align*}

% PCT-SOURCE: pdf=74 print=62 kind=prose
Thus, by (2-80),

% PCT-SOURCE: pdf=74 print=62 kind=display id=2-82
\begin{equation}
  |h^{(j)}(t)|\leq C\sup_\xi
  \frac{|P_\eta(\xi)|
  \left|\left(i\eta\mathbin{\cdot}\frac{\partial}{\partial\xi}\right)^j
  f(\xi)\right|(1+|\xi|^2)^k}{t^r}
  \tag{2-82}\label{eq:ch2-laplace-tempered-derivative-bound}
\end{equation}

% PCT-SOURCE: pdf=74 print=62 kind=prose
for some fixed sufficiently large $k$.  Now

% PCT-SOURCE: pdf=74 print=62 kind=display
\begin{equation*}
  h^{(j)}(t)=-\int_t^1d\tau\,h^{(j+1)}(\tau)+h^{(j)}(1)
\end{equation*}

% PCT-SOURCE: pdf=74 print=62 kind=prose
so, for $r>1$,

% PCT-SOURCE: pdf=74 print=62 kind=display id=2-83
\begin{equation}
  |h^{(j)}(t)|\leq
  \left(\frac{1}{t^{r-1}}-1\right)\frac{E}{r-1}+|h^{(j)}(1)| .
  \tag{2-83}\label{eq:ch2-laplace-tempered-integrated-bound}
\end{equation}

% PCT-SOURCE: pdf=75 print=63 kind=prose
where $E$ stands for the numerator in (2-82).  This last equation shows
that $h^{(j)}(t)$ is bounded in $t$ by a sum of terms, each of which goes to
zero as $f\to0$ in $\mathcal{S}$, and of degree $-(r-1)$ in $t$.  We insert
this expression in the formula for $h^{(j-1)}(t)$ and get an analogous
statement with degree $-(r-1)$ replaced by $-(r-2)$.  If we begin the
process with a $j$ sufficiently large, say $j=r+1$, and generalize (2-83)
to cover the case in which $1/\tau$ has to be integrated, we see that
$h(t)$ is bounded in $t$ by a sum of terms which approach zero as
$f\to0$ in $\mathcal{S}$.  Therefore

% PCT-SOURCE: pdf=75 print=63 kind=display
\begin{equation*}
  \lim_{t\to0}\int F(\xi-it\eta)f(\xi)\,d\xi
\end{equation*}

% PCT-SOURCE: pdf=75 print=63 kind=prose
exists and, by Theorems 2-6 and 2-9, $F$ is the Laplace transform of a
tempered distribution.
\end{proof}

% PCT-SOURCE: pdf=75 print=63 kind=prose
We can roughly summarize the results of Theorems 2-7 through 2-10: the
holomorphic functions which are Laplace transforms of tempered
distributions with supports in $\Gamma$ are holomorphic in the variable
$\xi-i\eta$ in $\mathcal{T}_n$, and of polynomial growth in $\eta$ near
infinity and zero.
